\documentclass[11pt]{beamer}
\usepackage[utf8]{inputenc}
\usepackage{amsmath, amssymb, graphicx, listings, xcolor}
\usepackage{bm}

% Presentation Theme
\usetheme{Madrid}
\usecolortheme{default}

% Code Block Styling
\lstset{
    language=Python, % Reusing Python highlighting for GMAT script readability
    basicstyle=\ttfamily\scriptsize,
    keywordstyle=\color{blue}\bfseries,
    commentstyle=\color{green!50!black},
    frame=single,
    backgroundcolor=\color{gray!10}
}

\title[Jupiter Polar Mission]{Astrodynamics 2026-1}
\subtitle{Advanced GMAT Targeting: \\ Polar Orbit Insertion \& Atmospheric Disposal}
\author{Prof. Juan}
\institute{Universidad de Antioquia - Aerospace Engineering}
\date{Spring 2026}

\begin{document}

% --- Title Slide ---
\begin{frame}
    \titlepage
\end{frame}

% ==========================================
% SECTION 1: MISSION ARCHITECTURE
% ==========================================
\section{Mission Architecture}

% --- Slide 1 ---
\begin{frame}{Mission Architecture: The Grand Finale}
    Our final mission simulation takes us from Low Earth Orbit (LEO) to a fiery disposal in the atmosphere of Jupiter. 
    
    \vspace{0.3cm}
    \textbf{The Three Distinct Phases:}
    \begin{enumerate}
        \item \textbf{Trans-Jovian Injection (TJI):} A patched-conic departure using a Differential Corrector to target the Jupiter arrival B-Plane.
        \item \textbf{Polar Orbit Insertion (JOI):} A massive retrograde braking burn at periapsis to capture the spacecraft into a $90^\circ$ inclination orbit around Jupiter.
        \item \textbf{Atmospheric Disposal:} A final apocenter maneuver to drop the periapsis deep into Jupiter's crushing atmosphere, ensuring planetary protection protocols.
    \end{enumerate}
\end{frame}

% ==========================================
% SECTION 2: B-PLANE TARGETING
% ==========================================
\section{Targeting the Polar Orbit}

% --- Slide 2 ---
\begin{frame}{The B-Plane: Aiming for the Poles}
    To achieve a polar orbit ($i = 90^\circ$), we cannot simply aim for the center of Jupiter. We must aim for a specific region of the \textbf{B-Plane}.
    
    \vspace{0.3cm}
    Recall that the B-Plane is perpendicular to our incoming hyperbolic excess velocity vector ($\bm{v}_\infty$).
    \begin{itemize}
        \item \textbf{$\bm{B} \cdot \bm{T}$ (Horizontal):} Controls the arrival node. For a polar orbit, we often aim straight down the center axis.
        \item \textbf{$\bm{B} \cdot \bm{R}$ (Vertical):} Controls the inclination. To fly exactly over the poles, our trajectory must be offset vertically in the B-Plane.
    \end{itemize}
    
    \vspace{0.2cm}
    \textit{In GMAT, our first Differential Corrector (DC1) will vary the Earth departure burn ($\Delta v$) to achieve the precise BdotT and BdotR needed for a $90^\circ$ arrival inclination.}
\end{frame}

% --- Slide 3 ---
\begin{frame}[fragile]{GMAT: The Arrival Targeting Sequence}
    Once we cross into Jupiter's Sphere of Influence, we use a second Differential Corrector (DC2) to execute the Jupiter Orbit Insertion (JOI).
    
\begin{lstlisting}
% 1. Target the Capture Burn at Periapsis
Target DC2 {SolveMode = Solve, ExitMode = SaveAndContinue};
    
    % Propagate to Periapsis (the most efficient place to burn)
    Propagate DefaultProp(Sat) {Sat.Jupiter.Periapsis};
    
    % Vary the Retrograde Burn
    Vary DC2(JOI_Burn.Element1 = -2.5); % Guess: 2.5 km/s braking
    
    Maneuver JOI_Burn(Sat);
    
    % Achieve the Desired Orbit 
    Achieve DC2(Sat.Jupiter.ECC = 0.85); % Highly elliptical capture
    Achieve DC2(Sat.Jupiter.INC = 90.0); % Polar orbit
    
EndTarget;
\end{lstlisting}
\end{frame}

% ==========================================
% SECTION 3: ATMOSPHERIC DISPOSAL
% ==========================================
\section{End of Mission: Atmospheric Disposal}

% --- Slide 4 ---
\begin{frame}{End of Mission: Planetary Protection}
    To prevent future contamination of potentially habitable moons (like Europa), mission profiles dictate that the spacecraft must be destroyed in Jupiter's atmosphere.
    
    \vspace{0.3cm}
    \textbf{The Orbital Mechanics of Disposal:}
    To lower the altitude of periapsis ($r_p$), we must decrease the orbital energy. The most efficient place to perform this maneuver is at \textbf{apoapsis} ($r_a$), where the spacecraft's velocity is lowest.
    
    \begin{equation}
        r_p = a(1-e) \quad \text{and} \quad r_a = a(1+e)
    \end{equation}
    
    By firing our engines retrograde (against the direction of motion) at apoapsis, we flatten the ellipse, driving the new periapsis below the 1-bar pressure level of Jupiter ($R_J \approx 71,492 \text{ km}$).
\end{frame}

% --- Slide 5 ---
\begin{frame}[fragile]{GMAT: The Disposal Sequence}
    Our final GMAT sequence utilizes a simple targeter to calculate the exact $\Delta v$ required to drop the periapsis into the atmosphere.
    
\begin{lstlisting}
% Propagate to the highest point in the orbit
Propagate DefaultProp(Sat) {Sat.Jupiter.Apoapsis};

Target DC3 {SolveMode = Solve, ExitMode = SaveAndContinue};
    
    % Vary the final retrograde maneuver
    Vary DC3(Death_Dive_Burn.Element1 = -0.5); 
    
    Maneuver Death_Dive_Burn(Sat);
    
    % Achieve a periapsis radius deep inside Jupiter
    % Jupiter Radius is ~71,492 km. Target 70,000 km.
    Achieve DC3(Sat.Jupiter.RMAG = 70000.0); 
    
EndTarget;

% Propagate the final plunge
Propagate DefaultProp(Sat) {Sat.Jupiter.RMAG = 70000.0};
\end{lstlisting}
\end{frame}

% --- Slide 6 ---
\begin{frame}{Summary}
    \textbf{Key Takeaways for Mission Design:}
    \begin{itemize}
        \item \textbf{Patched Conics $\rightarrow$ DC:} Analytical math gives us the initial guesses; GMAT's Differential Corrector makes it work in a perturbed 3D solar system.
        \item \textbf{B-Plane Mastery:} Controlling inclination requires precise targeting of the B-Plane before you even arrive at the planet.
        \item \textbf{Efficiency:} Always perform major orbit alterations at the apsides. Capture burns at periapsis (Oberth effect), and periapsis lowering burns at apoapsis.
    \end{itemize}
\end{frame}

\end{document}